% =========================================================
\section{Gas-Side Thermal Model}
\label{sec:thermal_model}
% =========================================================

\subsection{Model scope}

The implemented thermal module reports the local adiabatic-wall temperature,
\(T_{aw}(x)\), and a Bartz gas-side heat-transfer coefficient, \(h_g(x)\). The
coefficient is retained as a diagnostic indicator of where gas-side convective
coupling would be strongest; it is not combined with an independently prescribed
wall temperature. Thermophysical properties are supplied by RocketCEA and mapped
onto the quasi-one-dimensional flow reconstruction.

No conjugate wall-conduction, coolant or ablative model is solved. The adopted
gas-side boundary condition is therefore adiabatic,
\begin{equation}
\boxed{T_w(x)=T_{aw}(x)},
\qquad
\boxed{q''_w(x)=0}.
\label{eq:thermal_adiabatic_wall}
\end{equation}
Consequently, the thermal output must not be interpreted as a prediction of material
temperature or cooling-system heat load.

\subsection{Adiabatic wall temperature and recovery factor}

For the fully turbulent closure, the local recovery factor is derived from the local
edge Prandtl number,
\begin{equation}
r_f(x)=Pr_e(x)^{1/3}.
\label{eq:recovery_factor}
\end{equation}
The edge Prandtl number is obtained by interpolation of the RocketCEA chamber,
throat and exit transport values; it is not a user input. The adiabatic-wall
temperature is then
\begin{equation}
\boxed{
T_{aw}(x)=T_e(x)
\left[
1+r_f(x)\frac{\gamma_e(x)-1}{2}M_e(x)^2
\right]
}.
\label{eq:adiabatic_wall_temperature}
\end{equation}
This recovery-temperature condition is also used by the compressible
boundary-layer closures. The associated adiabatic Walz temperature--velocity
relation is developed in Sections~\ref{sec:integral_boundary_layer}
and~\ref{sec:blimp_cebeci_smith} \cite{Walz}.

\subsection{Diagnostic Bartz heat-transfer coefficient}

The gas-side convective heat-transfer coefficient is estimated using the simplified
Bartz correlation \cite{Bartz},
\begin{equation}
\begin{split}
h_g(x)={}&
\frac{0.026}{D_t^{0.2}}
\frac{\mu_t^{0.2}c_{p,t}}{Pr_t^{0.6}}
\left(\frac{p_c}{c^*}\right)^{0.8}
\left(\frac{D_t}{R_{c,t}}\right)^{0.1}
\left(\frac{A_t}{A(x)}\right)^{0.9}
\sigma(x),
\end{split}
\label{eq:bartz_hg}
\end{equation}
where \(D_t=2r_t\), \(R_{c,t}=0.4r_t\) is the divergent throat-arc
radius, and \(\mu_t\), \(c_{p,t}\) and \(Pr_t\) are evaluated at the
throat. The local area ratio follows directly from the generated contour,
\begin{equation}
\frac{A_t}{A(x)}=\left(\frac{r_t}{r(x)}\right)^2.
\end{equation}

The Bartz correction factor is evaluated consistently with the adiabatic wall,
\begin{equation}
\begin{split}
\sigma(x)={}&
\left[
\frac{1}{2}
\frac{T_{aw}(x)}{T_0}
\left(1+\frac{\gamma_e(x)-1}{2}M_e(x)^2\right)
+\frac{1}{2}
\right]^{-0.68}
\\
&\times
\left(1+\frac{\gamma_e(x)-1}{2}M_e(x)^2\right)^{-0.12}.
\end{split}
\label{eq:bartz_sigma}
\end{equation}
Although \(h_g\) generally peaks near the throat and remains useful for comparing
the spatial intensity of convective coupling, the present adiabatic boundary
condition gives
\begin{equation}
q''_g(x)=h_g(x)\left[T_{aw}(x)-T_w(x)\right]=0.
\label{eq:convective_heat_flux}
\end{equation}

\subsection{Interpretation of the adiabatic result}

Setting \(T_w=T_{aw}\) closes the gas-side boundary-layer problem without inventing
a geometry-independent solid temperature. It does not imply that a real nozzle is
thermally insulated, nor does it predict that the physical wall reaches the reported
temperature. A non-zero heat flux requires an additional thermal boundary condition
and a model that determines \(T_w\), for example conduction through the nozzle wall
coupled to regenerative, film or ablative cooling.

The current model omits wall thickness and material properties, transient energy
storage, radiation, coolant-channel geometry and coolant-side heat transfer. It
therefore supports aerodynamic preliminary design only. Thermal sizing and material
selection require a conjugate thermal model or experimental data after the geometry
has been defined.
